\section{The graphical interface}\label{sec:gui}

Launch with \code{uv run pyrate-ta-gui} (or \code{python -m pyrate_ta.gui}). The
window reuses
PyMORGAN's plot widget, plot-controls panel and theme so the two applications
look identical side by side.

\subsection{Layout}

The window is a single narrow control column on the left and a plot area on the
right; nothing is scattered across the width, so the default geometry stays
compact.

\begin{description}
  \item[Left column] Four stacked groups: \emph{Data} (load, clear and two
    lamps --- one for the dataset, one for the per-point noise; there is no
    format selector, since a fit needs a PDAT and its \code{.pdatn} sibling is
    picked up automatically), \emph{Model Selection}
    (\emph{Parallel} / \emph{Sequential} / \emph{Target} / \emph{ODE-defined}),
    \emph{Components} (the lifetime table with its bounds and fixed flags, and
    the add/remove buttons), \emph{IRF and time-zero}, and \emph{Fit options}
    (algorithm, restrict-to-display, noise weighting). The algorithm defaults
    to Trust Region Reflective: lifetimes are bounded below by zero, and
    MINPACK's Levenberg-Marquardt cannot honour bounds at all, so choosing it
    with finite bounds is refused rather than silently ignoring them.
  \item[Fit control] At the foot of the left column: the data / fit / residuals
    view selector, \emph{Preview}, \emph{FIT DATA}, \emph{Reset}, the statistic
    read-out, \emph{Copy taus}, and the two post-fit figures ---
    \emph{Plot DAS} (named after the model: DAS, EAS or SAS) and
    \emph{Plot Conc. Profile}. Both open in their own window, and both are
    opened automatically when a fit finishes; each has its own switch
    (\code{plot\_species\_after\_fit}, \code{plot\_profiles\_after\_fit}) in the
    \code{[gui]} section of \code{settings.toml}. A preview does not trigger
    them: it is a look at a guess, not a result.
  \item[Plot area] One figure, not three, so the window carries a single
    navigation toolbar. Its axes form a $3\times3$ grid: the contour takes the
    top-left $2\times2$ block (twice the size of the others), the kinetics
    panel the $2\times1$ column beside it and the transient spectra the
    $1\times2$ row below. The bottom-right cell holds no axis --- the
    \emph{Additional Plots and cuts} group sits there.

    The kinetics panel is drawn \emph{rotated} (PyMORGAN's
    \code{swap\_axes}), so its delay axis runs vertically like the contour's.
    Both shared axes then line up --- probe between the contour and the
    spectra, delay between the contour and the kinetics --- and the duplicated
    tick labels are dropped. Measured data are drawn as translucent points and
    the fit as a thin line over them, both configurable in the \code{[plots]}
    section of \code{settings.toml}.

    The legends of these panels are always drawn \emph{inside} their axes,
    whatever \code{[plot] legend\_location} says in PyMORGAN's settings: three
    panels share one figure, so a legend parked beside an axis lands on its
    neighbour or off the canvas. The setting still governs every stand-alone
    figure PyRATE-TA opens --- pop-out cuts, species spectra, concentration
    profiles --- where beside the axis is PyMORGAN's default and usual look.
  \item[Bottom] The plot-controls panel: axis limits, colour scale, contour
    levels, time-axis convention and probe units.
\end{description}

The embedded figure sets its margins explicitly (\code{[plots] panel\_left},
\code{panel\_right}, \code{panel\_top}, \code{panel\_bottom} and the two gaps)
rather than calling \code{tight\_layout}: it shares a canvas with the controls
panel, and the automatic layout leaves the axes noticeably smaller than the
window allows.

The four \emph{Additional plots and cuts} buttons open stand-alone PyMORGAN
figures. The two cut plots ask which cuts to draw, pre-filled with the
crosshair's position: the crosshair picks one cut, which is what the embedded
panels show, whereas a separate figure is where several are compared. The
prompt is PyMORGAN's (\code{gui.dialogs.ask\_values}), so the syntax ---
numbers, ranges such as \code{1900:5:1950}, or \code{all} --- is the same as in
PyMORGAN's own cut plots.

\subsection{Fitting on a small screen}

The window opens at $940 \times 900$ and is clipped at start-up to the
\emph{available} screen area --- which already excludes the menu bar and the
dock --- then centred, so a geometry saved on a large monitor cannot reopen
half under the dock on a laptop.

Being able to \emph{ask} for a small window is not enough: a layout whose
minimum size hint exceeds the screen simply refuses to shrink. The left column
therefore lives in a scroll area, so the window's minimum height is set by the
plot canvas alone and the controls stay reachable by scrolling when there is not
room for all of them at once.

Buttons are coloured by function, following the same palette as PyMORGAN --- a
contour button looks like a contour button in both applications --- and the
mapping lives in \code{gui/mw\_common.py} rather than in the \code{.ui}, so it
stays theme-responsive. Controls whose feature is not implemented yet are
disabled and carry the reason in their tooltip, rather than being present and
silently doing nothing.

\subsection{The K-matrix editor}

\emph{Define model...} opens the scheme editor: the scheme as text on the left
(section~\ref{sec:notation}), its graph and the derived rate matrix on the
right. \emph{Update K graph} redraws \emph{and} checks --- a scheme that does
not parse produces its error, with the line number, instead of a picture, so
the button answers ``would this scheme work?'' before a fit is started, and
\emph{OK} stays disabled until it does. The matrix is shown with every rate set
to $1$, since what is being checked there is the structure, not the values.
Accepting the scheme sets the lifetime table to one row per rate constant,
named after it.

\emph{Show kinetic graph}, beside it, draws the scheme of whatever model is
currently selected --- compartments, arrows and rate \emph{symbols}
($k_1$, $k_2$, \dots, not their values) --- in its own window. It works before
any fit, from the lifetimes as typed, so a scheme can be checked while it is
being built; after a fit it draws the fitted one.

\subsection{The LDA tab}

A separate \emph{LDA} tab gives access to Lifetime Density Analysis
(section~\ref{sec:lda}) without switching away from the same dataset. Its
controls map directly to the keyword arguments of \code{solve\_lda}:

\begin{description}
  \item[Grid] Number of log-spaced lifetime points and the
    $[\tau_\text{min},\,\tau_\text{max}]$ range (auto-filled from the delay axis).
  \item[Regularisation] Penalty type (\code{ridge} / \code{d1} / \code{d2}) and
    the alpha method (\code{lcurve} / \code{gcv} / \code{morozov} / \code{fixed}).
    When \emph{fixed} is selected an alpha entry box appears.
  \item[Options] Non-negative constraint, SVD pre-filtering (number of
    components), coherent-artefact suppression, and bootstrap iterations.
  \item[Run LDA] Executes \code{solve\_lda} headlessly and populates the
    results table and the lifetime-density map.
\end{description}

After a successful run the tab shows: the integrated amplitude spectrum
$A(\tau) = \sum_p |S_{p}(\tau)|$, the L-curve or GCV score plot (for diagnosing
the alpha choice), and a 2D heatmap $S(\tau,\lambda)$. Detected peaks are
listed in a table with their centroid lifetimes. The result can be exported to a
\code{.pdat} file with \emph{Export map\ldots} for further plotting in PyMORGAN.

\subsection{Reading a fit}

Two figures carry the result, and both are opened automatically when the fit
converges (buttons reopen them at will):

\begin{description}
  \item[Species spectra] The amplitude spectrum of each component, drawn by
    PyMORGAN's \code{plot\_species\_spectra} through the result's
    \code{as\_species\_args()}. Their name follows the model, and so does the
    button: DAS for parallel, EAS for sequential, SAS for target. The
    distinction is not cosmetic --- it is what the spectra may be interpreted
    as.
  \item[Concentration profiles] The populations $C(t)$ of the model itself: one
    of the few views PyMORGAN has no concept of, since they are never measured.
    Component $i$ carries the same colour as its spectrum, and the delay axis
    comes from PyMORGAN's \code{apply\_delay\_axis}, so the profiles can be laid
    beside the contour and read on the same time axis. For a sequential or
    target model this is where the plot earns its keep: it shows \emph{when}
    each species is present, and therefore which spectrum can be trusted where.
\end{description}

\subsection{Layout is declared, not built}

All widgets, layouts, containers and controls are declared in the Qt Designer
file \code{src/pyrate_ta/gui/main\_window.ui}. Python code binds signals, slots,
dynamic styles and data states to widgets defined there; it does not construct
them. Edit the layout with

\begin{lstlisting}[language=bash, style=py]
uv run pyrate-ta-edit-gui        # or: python -m pyrate_ta.gui.designer
\end{lstlisting}

The plot area is promoted to PyMORGAN's \code{WidgetPlot}, and the whole
plot-controls group (\code{PC\_box}) is the one from PyMORGAN's layout, name for
name, so that \code{pymorgan.gui.plot\_controls.PlotControlsPanel} drives it
unmodified. Renaming a \code{PC\_*} widget breaks that contract; a test asserts
the full set is present.

\subsection{Structure of the code}

\code{MainWindow} is assembled from per-tab mixins in \code{gui/tabs/}:

\begin{description}
  \item[\code{DataTabMixin}] Loading, dataset display, the embedded
    data / fit / residuals view and the additional cut plots.
  \item[\code{ModelTabMixin}] Model selection, the lifetime table, bounds,
    fixed flags, IRF and time-zero controls.
  \item[\code{FittingTabMixin}] Fit control: run, preview, reset, statistic
    read-out, copy taus, auto-plots.
  \item[\code{LDATabMixin}] The LDA controls and result display.
\end{description}

A mixin is a plain method container with no state of its own, so \code{self} is
always the main window. \code{main\_window.py} keeps construction, wiring,
menus and theming only; shared constants live in \code{gui/mw\_common.py}.

Wiring is defensive: a connection is made only if the named widget exists, so
the window still loads while the layout is being edited in Designer, and the
missing widget is logged rather than raising at start-up.

\subsection{Rendering}

A contour re-plot cannot be updated artist by artist --- the level set itself
changes --- so rapid control changes, such as a colour-scale slider drag that
emits one signal per step, are coalesced into a single render through a
short-interval timer. Limit changes that do not alter the levels re-apply the
axis limits without re-plotting.

\subsection{Start-up cost}

\code{gui/app.py} imports neither \code{pyrate} nor \code{pymorgan} at module
level, so the window (and, where enabled, the splash screen) can appear before
scipy, matplotlib and the loader registries are paid for. The same discipline
applies to the package itself: \code{pyrate/\_\_init\_\_.py} binds its exports
through a PEP~562 \code{\_\_getattr\_\_}, so \code{import pyrate_ta} is nearly
free and each module is imported on first use of a name it provides.
